% s-254-p254-m002; original PDF 254, 255
\Aphorism{s-254-p254-m002}{351}
Ein Künstler kann und wird in der Regel dem, was er als Künstler schafft, als Mensch so ahnungslos und unfähig gegenüberstehen wie der Erstbeste aus dem Publicum; das ist ein altes und lustiges Paradoxon. Fremde schmücken und taufen das Kind; die Mutter liegt im Bette und schläft sich aus. Der Künstler, dessen Werk unter Umständen eine neue Cultur und Gesittung heraufführt, ist mit seiner Moral, seiner Bildung ganz und gar rückständig und in erhabener Weise zum Missverständniss seiner selbst geneigt. Der und der Schauspieler strahlt eine nahezu metaphysische Wirkung aus, er redet wie aus dem „Abgrunde“, wie Einer, der eine unmögliche Erfahrung hinter sich hat, der in einer entsetzlichen Einsamkeit und Unterwelt einmal zu Gaste gewesen ist — und dieser geborene Prinz Hamlet dachte mit seinem Oberflächenbewusstsein vielleicht an sein Honorar oder an die Zeitungskritik, als er dergestalt als Künstler auf seine Höhe kam. Mozarts Don Juan: nach des guten alten Meisters Absicht unzweifelhaft ein demoralisirtes Ungeheuer, das \Foreign{a priori} Unrecht hat und an dem mit Furcht, Mitleid und tugendhafter Entrüstung nicht gespart werden soll — für uns etwas unvergleichlich Höheres: ein Märtyrer des Genusses, dem es hinsichtlich des Weibes nicht beschieden ist bei einer einzelnen Enttäuschung stehen zu bleiben, eine skeptische selbstquälerische verachtende Seele mit einem ungeheuren Zug ins Tragische und Fatalistische, ein unbegriffener unbegreiflicher Grandseigneur, stolz und ironisch genug, um den harmlosen Moralweiblein, die er auf dem Gewissen hat, noch dadurch Recht zu geben, dass er sich von ihrem Teufel holen lässt. So versteht Francesco d’Andrade den Don Juan; welcher deutsche Darsteller wäre vorurtheilsfrei genug, ihn ebenso zu verstehen — also vor allen Dingen unmozartisch! Denn Mozart wusste nicht, was er that, und vielleicht kein Künstler weiss es — diese furchtbaren Unbewussten, die in allerunterthänigster Beschränktheit ihre zeitgemässe Cultur nur ein wenig beleuchten, erleuchten wollen und dabei das ganze Zeitalter in Brand setzen!
\par

